\chapter{Getting Started}

Das Repository \textbf{loclass-starter} ist bewußt übersichtlich aufgebaut. Ziel ist es, einen möglichst einfachen Einstieg in das Framework zu ermöglichen, ohne dass zunächst Änderungen an der Dokumentklasse notwendig sind.

\section{Projekt klonen}

Zunächst wird das Repository mit Git in ein beliebiges Verzeichnis geklont.

\begin{quote}
\texttt{git clone <Repository-URL>}
\end{quote}

\section{Dokument erstellen}

Zum Erzeugen des Dokuments genügt der Aufruf

\begin{quote}
\texttt{latexmk}
\end{quote}

Der Build-Prozess erzeugt automatisch alle erforderlichen Hilfsdateien und erstellt anschließend die PDF-Datei \texttt{main.pdf}.

Weitere Schritte sind nicht erforderlich.

\section{Projektstruktur}

Für die tägliche Arbeit sind im Wesentlichen drei Verzeichnisse relevant.

\subsection*{project}

Dieses Verzeichnis enthält die projektspezifische Konfiguration, beispielsweise Metadaten, eigene Makros oder zusätzliche Pakete.

\subsection*{content}

Dieses Verzeichnis enthält den eigentlichen Inhalt des Dokuments.

Die Kapitel werden durch einen zweistelligen Präfix automatisch in die gewünschte Reihenfolge gebracht. Beispielsweise:

\begin{itemize}
    \item \path{10_introduction.tex}
    \item \path{20_getting_started.tex}
    \item \path{30_project_structure.tex}
\end{itemize}

\begin{quote}
Die Dateien werden beim Build automatisch erkannt und in numerischer Reihenfolge in das Dokument eingebunden.
\end{quote}

\subsection*{assets}

Dieses Verzeichnis enthält statische Ressourcen wie Bilder, Logos oder weitere Dokumente, die im Dokument verwendet werden.
